\subsection{The politico-economic equilibrium with log}\label{\refsec:PEELOG}
With log preferences the politico-economic equilibrium is inexpensive to solve, and the reason why is worth stating before the algorithms, because it is the structure the algorithms exploit. Let $z_t$ denote the marginal effect of the bounded policy on the political objective, for each $t$ defined as
\begin{align}\label{\refeq:LOG:fast}
    z_t &= \omega\sum_{j}\left(\gamma_{t-1,j}p_{t-1,j}\mu_{t-1,j} \frac{d\upsilon_{2,t}^j}{d\tau_t}\right)+\nu_t\sum\left(\gamma_{t,j}\mu_{t,j} \frac{d\upsilon_{1,t}^j}{d\tau_t}\right),
\end{align}
where the marginal effects on indirect utilities are defined by equations \eqref{\refmodeleq:PEELOG} for $t<T$ and \eqref{\refmodeleq:terminalPEELOG} for $t=T$.

\smalltitle{The problem carries no endogenous state, and the system is triangular.}
Two structural facts shape everything below. The first is the state-space reduction of \cref{\refmodelsec:PEE}: the savings shares $s_{t-1,i}/s_{t-1}$ that the old generations carry into period $t$ are closed-form functions of $\tau_t$ (evaluated at the previous period's parameter vintage), so the distribution of past savings never has to be tracked as a state, and under log preferences no other predetermined quantity enters $z_t$ either. Each $z_t$ is a function of current and future policy alone.

The second fact concerns \emph{which} future policy. Inspecting the marginal effects in \eqref{\refmodeleq:PEELOG}, the two young terms $d\upsilon_{1,t}^i/d\tau_t$ and $d\upsilon_{1,t}^0/d\tau_t$ depend on $\tau_t$ alone. The old formal term $d\upsilon_{2,t}^i/d\tau_t$ depends on $\tau_t$ directly and through $s_{t-1}^i/s_{t-1}$. The old informal term $d\upsilon_{2,t}^0/d\tau_t$ additionally involves the level of $\Theta_{h,t}$, which by \eqref{\refeq:auxiliary:Thetah} depends on $\tau_t$ and on $\tau_{t+1}$ through $\Gamma_{s,t}$. No term depends on any lag of $\tau$. Hence
\begin{align}\label{\refeq:recursive}
    z_t = z_t\left(\tau_t, \tau_{t+1}\right), \qquad\text{and}\qquad z_T = z_T(\tau_T),
\end{align}
the latter because the terminal period uses $\Theta_{h,T}(\tau_T)$ from \eqref{\refeq:auxiliary:ThetahT} in place of \eqref{\refeq:auxiliary:Thetah}.

The system is therefore triangular: it can be solved exactly by a \emph{backward recursion} in which each period is a \emph{scalar} problem in $\tau_t\in[l,u]$, given the already-solved $\tau_{t+1}$. This is what makes robust one-dimensional methods available at all, and it removes the dependence on an initial guess for the entire path. Note that the property relies on $\theta_t$ and $\epsilon_t$ being exogenous; if pension characteristics are chosen alongside $\tau_t$, the argument has to be revisited.

\smalltitle{Behaviour at the boundaries.}
Both $d\upsilon_{1,t}^j/d\tau_t$ and $\partial \ln(\Theta_{h,t})/\partial\tau_t$ carry a factor $1/(1-\tau_t)$, so $z_t\rightarrow -\infty$ as $\tau_t\rightarrow 1$. We therefore always evaluate on $[l,u]$ with $u<1$ strictly. The divergence makes it easy to choose $u$ close enough to $1$ that $z_t(u)<0$, and we assume this throughout. Two consequences follow: the upper corner is then never selected, and $z_t(l)>0$ guarantees by continuity that at least one interior downward crossing exists. The substantive difficulty is thus not existence, but (i) detecting a lower-corner solution and (ii) choosing between several interior crossings when they occur --- exactly the difficulties \cref{\refsec:RobustRoot} is built for. Both are verified numerically rather than assumed, since $z_t(u)<0$ is a property of the calibration, not an identity.

\begin{alg}[Gradient-based roots for the politico-economic equilibrium with log preferences]
\label{\refalg:LOG:fast}
This is barely an algorithm, in the sense that it includes no fallback options, but simply presumes that the first order condition is well behaved. Let $\tilde{\bm{\tau}} \in \mathbb{R}^{T}$ be a candidate vector of policies and $\bm{\tau}$ the bounded version where each $\tau_t \in [l, u]$. Stack \eqref{\refeq:root} across periods, with $z_t$ from \eqref{\refeq:LOG:fast}, and solve for the $\tilde{\bm{\tau}}$ that zeroes the stacked system with a standard gradient-based root finder; report $\bm{\tau}$ as the bounded version. By \eqref{\refeq:recursive} the system's Jacobian is upper bidiagonal, but the solver is not told so; the value of this variant is speed and machine precision when it works, not robustness.
\end{alg}

\begin{alg}[Backward grid search]\label{\refalg:LOG:gridsearch}
Fix $\bm{\theta},\bm{\epsilon}$ and a global grid $\mathcal{T}$, extended to $\tilde{\mathcal{T}}$ as in \eqref{\refeq:extendedGrid}. Exploiting \eqref{\refeq:recursive}, solve backwards.

\smallskip\noindent\emph{Terminal period.} Evaluate $z_T$ on all of $\mathcal{T}$. \emph{Initial boundary check}: if $z_T<0$ throughout, set $\tau_T = l$; if $z_T>0$ throughout, set $\tau_T = u$. Otherwise construct the candidate set $\mathcal{C}_T$ from \eqref{\refeq:candidates} and set $\tau_T = \arg\max_{\mathcal{C}_T}\hat{\mathcal{V}}_T$.

\smallskip\noindent\emph{Recursion, $t = T-1,...,1$.} Given the solved $\tau_{t+1}$:
\begin{enumerate}
    \item \emph{Refine.} Let $\mathcal{T}_t\subseteq\mathcal{T}$ consist of the nodes from $\Delta^-$ below to $\Delta^+$ above the node nearest $\tau_{t+1}$. The window need not be symmetric: where the policy path is known to be monotone, $\Delta^-\neq\Delta^+$ saves evaluations.
    \item \emph{Evaluate and select.} Evaluate $z_t(\cdot,\tau_{t+1})$ on $\mathcal{T}_t$ and apply the terminal-period rule to obtain a candidate $\tau_t$.
    \item \emph{Expand.} If the selected $\tau_t$ falls on an endpoint of $\mathcal{T}_t$ that is not also an endpoint of $\mathcal{T}$, the window may have excluded the maximum. Widen $\mathcal{T}_t$ in that direction and return to step 2, re-using the nodes already evaluated.
\end{enumerate}
The cost is $O(T\cdot M)$ evaluations of closed-form expressions, with no derivatives and no initial guess. The accuracy of the located interior roots is capped by the grid resolution.
\end{alg}

\smalltitle{The recommended default combines the two.}
Neither algorithm alone is the production solver. The practical strategy is: attempt Algorithm \ref{\refalg:LOG:fast} first; if its residual fails the convergence check, run Algorithm \ref{\refalg:LOG:gridsearch} and polish the resulting path with Algorithm \ref{\refalg:LOG:fast}, warm-started at the grid solution; if the polish itself fails, return the grid solution flagged as grid-accurate. The division of labour is deliberate: the grid search locates the correct branch of a problem with corners and possible multiplicity, and the gradient step then attains machine precision on a problem that, near the located solution, has neither.
